% 自动生成：gen_refs.py（参考文献照录原书，未翻译；自动转录，未逐条校对）
\chapter*{参考文献}
\addcontentsline{toc}{chapter}{参考文献（原书 p.385–404）}
\markboth{参考文献}{参考文献}

{\small
\noindent （原则上只收录正文引用过的文献，并不意图网罗辛几何的重要文献。文献末尾所附 hep-th/…、math/… 等为互联网论文服务器 eprint 的参照编号，这些文献可从京都大学基础物理学研究所的镜像服务器 http://www.yukawa.kyoto-u.ac.jp/ 下载。）\par\medskip

\noindent [1]\ 足立正久，『微分位相幾何学』（現代の数学14), 共立出版， 1976.\par\smallskip
\noindent [2]\ 足立正久，『埋め込みとはめ込み』（数学選書），岩波書店， 1984.\par\smallskip
\noindent [3]\ Adams, J., Infinite Loop Spaces, Annals of Math. Studies 90, Princeton Univ. Press, 1978.\par\smallskip
\noindent [4]\ Agrachev, A., "Method of control theory in nonholonomic geometry" in International Congress of Mathematics, Birkhauser, Zurich, 1995.\par\smallskip
\noindent [5]\ Akbult, S. and McCarthy, J., Casson Invariant for Oriented Homology 3-Spheres, an Exposition, Mathematical Notes 36, Princeton Univ. Press, 1990.\par\smallskip
\noindent [6]\ Andersen, J., Mattes, J. and Reshetikhin, N., The Poisson structure on the moduli space of flat connections and chord diagrams, Topology 35(1998), pp.1069--1083.\par\smallskip
\noindent [7]\ Andersen, J., Mattes, J . and Reshetikhin, N., Quantization of the algebra of chord diagrams, Math. Proc. Cambridge Philos. Soc. 124(1998), pp.451-463.\par\smallskip
\noindent [8]\ Arnold, V., On the representation of a continuous function of three variables by superpositions of functions of two variables, Math. Sbomik 48(1959), pp. 3-74. (in Russian)\par\smallskip
\noindent [9]\ Arnold, V., Sur une proprietes des application globalement canoniques de la mecanique classique, C. R. Acad. Sci. Paris 261(1965), pp. 3719-3722.\par\smallskip
\noindent [10]\ Arnold, V., "Mathematical developments arising from Hilbert problems" in Proc. Sympos. PuだMath., 1966, p. 66.\par\smallskip
\noindent [11]\ Arnold, V., Normal forms for functions near degenerate critical points, the Weyl group Ak, Dk, Ek and Lagrangian singularities, Functional Anal. Appl. 6(1972), pp. 254-272.\par\smallskip
\noindent [12]\ Arnold, V., Lagrange and Legendrian cobordism I , II, Functional Anal. Appl. 14(1980), pp.167-177, 252-260.\par\smallskip
\noindent [13]\ Arnold, V., Mathematical Methods of Classical Mechanics, Graduate Texts in Math. 60, Springer, Berlin, 1989. (邦訳）『古典力学の数学的方法』，安藤詔ー·蟹江 幸博·丹羽敏雄訳，岩波書店， 1980-\par\smallskip
\noindent [14]\ Arnold, V., Singularity of Caustics and Wave Fronts: Mathematics and its applications, Kulwer Academic Publishers, 1990.\par\smallskip
\noindent [15]\ Arnold, V. and Givental, A., "Symplectic geometry" in Encyclopaedia of Mathematical Sciences, Dynamical system IV, Arnold, V. and Novikov, S.(eds.), Springer, Berlin, 1980, pp. 4-136.\par\smallskip
\noindent [16]\ Arnold, V., Guestein, S. and Varchenko, A., Singularities of Differential Maps I , II, Monographs in Mathematics 82, 83, Birkhauser, Basel, 1985, 1988.\par\smallskip
\noindent [17]\ Arnold, V. and Khesin, B., Topological Methods in Hydrodynamics, Applied Mathematical Sciences 125, Springer, Berlin, 1998.\par\smallskip
\noindent [18]\ Atiyah, M., Convexity and commuting Hamiltonians, Bulletin of the London Math. Soc. 14(1982), pp.1-15.\par\smallskip
\noindent [19]\ Atiyah, M. and Bott, R., The Yang-Mills equations over Riemann surfaces, Philosophical Trans. of the Royal Soc. of London Ser. A, 308(1982), pp. 523-615.\par\smallskip
\noindent [20]\ Atiyah, M. and Bott, R., The moment map and equivariant cohomology, Topology 23(1982), pp.1-28.\par\smallskip
\noindent [21]\ Atiyah, M. and Hitchin, N., The Geometry and Dynamics of Magnetic Monopoles, Princeton Mathematical Series 36, Princeton Univ. Press, 1989.\par\smallskip
\noindent [22]\ Atiyah, M., Patodi, V. and Singer, I . , Spectral asymmetry and Riemannian geometry I , II, III, Math. Proc. Camb疇e Philos. Soc. 77, 78, 79(1975), pp. 43-69, 405-432, 71-99.\par\smallskip
\noindent [23]\ Audin, M., Quelques calculs en cobordisme lagrangien, Ann. Inst. Fourier 35(1985), pp.159-194.\par\smallskip
\noindent [24]\ Audin, M., The Topology of Torus Actions on Symplectic Manifolds, Progress in Math. 93, Birkhauser, Basel, 1991.\par\smallskip
\noindent [25]\ Audin, M. and Lafontaine, J.(eds.), Holomorphic Curves in Symplectic Geometry, Progress in Math. 117, Birkhauser, Basel, 1994.\par\smallskip
\noindent [26]\ Austin, D. and Braam, P., "Morse-Bott theory and equivariant cohomology" in The Floer Memorial Volume, Hofer, H. et al.(eds.), Birkhauser, 1995, pp. 123-184.\par\smallskip
\noindent [27]\ Axelrod, S. and Singer, I . , "Chern-Simons perturbation theory I" in Proc. XXth International Conference on Differential Geometric Method in Theoretical Physics, Catto, S. and Rocha, A.(eds.), World Scientific, Singapore, 1991, pp.3-45.\par\smallskip
\noindent [28]\ Axelrod, S. and Singer, I . , Chern-Simons perturbation theory II, J. Di.ff. Geom. 39(1991), pp. 173-213\par\smallskip
\noindent [29]\ Bando, S., Kasue, A. and Nakajima, H., On a construction of coordinates at infinity on manifolds with fast curvature decay and maximal volume growth, Invent. Math. 97(1989), pp. 313-349.\par\smallskip
\noindent [30]\ Banks, T., Matrix theory, Nucl. Phys. Proc. Suppl. 67(1998), pp.180-224. hepth/9710231.\par\smallskip
\noindent [31]\ Banyaga, A., The Structure of Classical Diffeomorphism Groups, Mathematics and its Applications 400, Kluwer Academic Publishers, 1997.\par\smallskip
\noindent [32]\ Barannikov, S., Generalized periods and mirror symmetry in dimension n > 3. math/9903124.\par\smallskip
\noindent [33]\ Barannikov, S. and Kontsevich, M., Frobenius manifolds and formality of Lie algebras of polyvector fields, Intemat. Math. Res. Notices 4(1998), pp. 201-215\par\smallskip
\noindent [34]\ Bates, S. and Weinstein, A., Lectures on the Geometry of Quantization, Berkeley Mathematics Lecture Notes 8, Amer. Math. Soc., 1991.\par\smallskip
\noindent [35]\ Bauendi, M., Ebenfelt, P. and Rothschild, L., Real Submanifolds in Complex Space and their Mappings, Princeton Mathematical Series 47, Princeton Univ. Press, 1999.\par\smallskip
\noindent [36]\ Bayen, F., Flato, M., Fronsdal, C., Lichnerowicz, A. and Sternheimer, D., Deformation theory and quantization I , II, Ann. of Phys. 111(1978), pp. 61-110, 111-151.\par\smallskip
\noindent [37]\ Becchi, C., Rouet, A. and Stora, R., Renormalization of gauge theories, Ann. of Phys. 98(1974), p. 287. .\par\smallskip
\noindent [38]\ Becker, K., Becker, M. and Strominger, A., Fivebranes, membranes and nonperturbative string theory, Nucl. Phys. B456(1995), pp.130-152. hep-th/9507158.\par\smallskip
\noindent [39]\ Behrend, K., Gromov-Witten invariants in algebraic geometry, Invent. Math. 127(1997), pp. 604--617.\par\smallskip
\noindent [40]\ Behrend, K. and Fantechi, B., The intrinsic normal cone, Invent. Math. 128(1997), pp. 45-88.\par\smallskip
\noindent [41]\ Bella'iche, A. and Risler, J.(eds.), Sub-Riemannian Geometry, Progress in Math. 144, Birkhiiuser, Basel, 1991.\par\smallskip
\noindent [42]\ Bennequin, D., Entracements et 恥ation de Pfaff, Asterisque 107-108, 1983, pp. 87-161.\par\smallskip
\noindent [43]\ Bershadsky, M., Ceccoti, S., Ooguri, S. and Vafa, C., Kodaira-Spencer theory of gravity and exact results for quantum string amplitudes, Comm. Math. Phys. 165(1994), pp. 311--427.\par\smallskip
\noindent [44]\ Bini, G., de Concini, S., Polito, M., Procesi, C., On the work of Givental relative to mirror symmetry, Appunti dei Corsi Tenuti da Docenti della Scuola, Scuola Normale Superiore, Pisa, 1998.\par\smallskip
\noindent [45]\ Boardman, J. and Vogt, R., Homotopy Invariant Algebraic Structures on Topological Spaces, Lecture Note in Math. 347, Springer, Berlin, 1973.\par\smallskip
\noindent [46]\ Bogomolov, F., Hamiltonian Kahler manifolds, Dokl. Akad. Nauk SSSR 243 (1978), pp. 1101-1104\par\smallskip
\noindent [47]\ Borcherds, R., Automorphic forms with singularities on Grassmannians, Invent. Math. 132(1998), pp. 491-562.\par\smallskip
\noindent [48]\ Borcea, C., "K3 Surfaces with Involution and Mirror Pairs of Calabi-Yau Manifolds" in Mirror Symmetry II, Greene, B. and Yau, S. T.(eds.), International Press, Hong-Kong, 1997, pp. 717-744.\par\smallskip
\noindent [49]\ Borel, A., Sur la cohomologie des espaces fibres principaux et des espaces homogらnes de groupes de Lie compacts, Ann. of Math. 57(1953), pp.115-207.\par\smallskip
\noindent [50]\ Bott, R., Nondegenerate critical manifolds, Ann. of Math. 60(1954), pp. 248--261.\par\smallskip
\noindent [51]\ Bourgain, J., "Harmonic analysis and nonlinear partial differential equations" in International Congress of Mathematics, Birkhiiuser, Ziirich, 1995.\par\smallskip
\noindent [52]\ Bryant, R., Metrics with exceptional holonomy, Ann. of Math. 126(1987), pp. 525-576.\par\smallskip
\noindent [53]\ Bryant, R., Chern, S., Gardner, R., Goldschmidt, H. and Griffiths, P., Exte-rzor Differential Systems, Mathematical Sciences Research Institute Publications 18, Springer, New York, 1991.\par\smallskip
\noindent [54]\ Burago, Y. and Zalgaller, A., Geometric Inequalities, Springer, Berlin, 1988.\par\smallskip
\noindent [55]\ Calabi, E., "On the group of automorphisms of a symplectic manifold" in Problems in Analysis, Gunning, C. (ed.), Princeton Univ. Press, 1970, pp.1-26.\par\smallskip
\noindent [56]\ Candelas,P., delaOssa, X. C., Green, P. S. and Parks, L., A pair of Calabi-Yau manifolds as an exactly soluble superconformal theory, Nucl. Phys. B359(1991), p. 21.\par\smallskip
\noindent [57]\ Carrion, R., A generalization of the notion of instanton, Diff. Geom. and its Application 8(1998), pp. 1-20:\par\smallskip
\noindent [58]\ Cattaneo, A. and Felder, G., The Path Integral Approach to the Kontsevich Quantization Formula. math.QA/9902090, Comm. Math. Phys. 212(2000), no. 3, pp. 591-611.\par\smallskip
\noindent [59]\ Chaperon, M., Quelques Questions de Geometrie Symplectique (d'aprもs, entre autres, Poincare, Arnol'd, Conley et Zehnder), Asterisque 105-106, 1983, pp. 231-249.\par\smallskip
\noindent [60]\ Chekanov, Differential algebras of Legendrian links, Invent. Math. 150(2002), no. 3, pp. 441-483.\par\smallskip
\noindent [61]\ Conley, C. and Zehnder, E., The Birkhoff-Lewis fixed point theorem and a conjecture by V. I . Arnold, Invent. Math. 73(1983), pp. 33-49.\par\smallskip
\noindent [62]\ Connes, A., Noncommutative Geometry, Academic Press, 1994.\par\smallskip
\noindent [63]\ Connes, A., Douglas, R. and Schwarz, A., Noncommutative geometry and matrix theory: compactification on tori, J. High Energy Phys 2(1998). Paper 3 hepth/9711162.\par\smallskip
\noindent [64]\ Cox, D. and Katz, S., Mirror Symmetry and Algebraic Geometry, Mathematical surveys and Monographs 68, Amer. Math. Soc., 1998.\par\smallskip
\noindent [65]\ Dazord, P. and Weinstein, A.(eds.), Symplectic Geometry, Groupoids, and Integrable Systems, Mathematical Sciences Research Institute Publications 20, Springer, New York, 1989.\par\smallskip
\noindent [66]\ Deligne, P., Etingof, P., Freed, D., Jeffrey, L., Kazdan, D., Morgan, J., Morrison, D. and Witten, E. (eds.), Quantum Fields and Strings: A course for mathematicians, , 2, Amer. Math. Soc., 1999.\par\smallskip
\noindent [67]\ Deligne, P. and Freed, D., "Supersolutions" in Quatum Fields and Strings: A course for mathematicians 1, 2, Deligne, P. et al.(eds.), Amer. Math. Soc., 1999.\par\smallskip
\noindent [68]\ Deligne, P., Griffiths, P., Morgan, J. and Sullivan, D., Real homotopy theory of Kahler manifolds, Invent. Math. 29(1975), pp. 245-274.\par\smallskip
\noindent [69]\ Deligne, P. and Mumford, D., The irreducibility of the space of curves of given genus, Publ. Math. IHES 36(1969), pp. 75-110.\par\smallskip
\noindent [70]\ De Wilde, M. and Lecomte, P., "Formal deformations of the Poisson Lie algebra of a symplectic manifold and star-products. Existence, equivalence, derivations." in Deformation Theory of Algebras and Structures and Applications, Hazewinkel, H. and Gerstenhaber, M.(eds.), Kluwer Academic Publishers, 1988.\par\smallskip
\noindent [71]\ Donagi, R., "Taniguchi lecture on principal bundles on elliptic fibrations", in Integrable Systems and Algebraic Geometry, World Scientific, 1998, pp. 33--46. hep— th/9802094. (\par\smallskip
\noindent [72]\ Donaldson, S., A new proof of a theorem of Narasimhan and Seshadri, J. Diff. Geom. 18(1983), pp. 279-315. (\par\smallskip
\noindent [73]\ Donaldson, S., Anti-self-dual Yang-Mills connections over complex algebraic surfaces and stable vector bundles, Proc. London Math. Soc. 50(1985), pp.1-26. (\par\smallskip
\noindent [74]\ Donaldson, S., Irrationality and h-cobordism conjecture, J. Diff. Geom. 26(1986), pp. 275-297. (\par\smallskip
\noindent [75]\ Donaldson, S., Infinite determinants, stable bundles and curvature, Duke Math. J. 54(1987), pp. 231-241. (\par\smallskip
\noindent [76]\ Donaldson, S., Lecture at University of Warwick, 1992. (\par\smallskip
\noindent [77]\ Donaldson, S., Symplectic submanifolds and almost-complex geometry, J . Diff. Geom. 44(1996), pp. 275-297. (\par\smallskip
\noindent [78]\ Donaldson, S., "Lefschetz fibrations in symplectic geometry" in International Congress of Mathematics, 1998, Berlin.\par\smallskip
\noindent [79]\ Donaldson, S. and Kronheimer, P., The Geometry of Four Manifolds, Oxford Univ. Press, 1990. (\par\smallskip
\noindent [80]\ Donaldson, S. and Thomas, R., "Gauge theory in higher dimensions" in The Geometric Universe, Oxford Univ. Press, 1996, pp. 31--47. (\par\smallskip
\noindent [81]\ Dubrovin, B., "The geometry of 2D topological field theories" in Integrable sys— tems and quantum groups, Springer, Berlin, 1996, pp.120-348. (\par\smallskip
\noindent [82]\ Dubrovin, B., Painleve trancendents and two-dimensional topological field the。ry. math/9803017. The Painleve property, pp. 287-412, CRM Ser. Math. Phys., Springer, New York, 1999. (\par\smallskip
\noindent [83]\ Dubrovin, B., Krichever, I . and Novikov, S., "Integrable systems I" in Encyclopaedia of Mathematical Sciences, Dynamical system IV, Arnold, V. and Novikov, S.(eds.), Springer, Berlin, 1980, pp.174-283.\par\smallskip
\noindent [84]\ Dubrovin, B. and Zhang, Y., Frobenius manifolds and Virasoro constraints. math/9808048. Selecta Math. (N.S.) 5(1999), no. 4, pp. 423-466. (\par\smallskip
\noindent [85]\ Duistermaat, J. and Heckman, G., On the variation in the cohomology class of the symplectic form of the reduced phase space, Invent. Math. 69(1982), pp. 259-269.\par\smallskip
\noindent [86]\ Eells, J . and Sampson, J., Harmonic mappings of Riemannian manifolds, Amer. J. Math. 86(1964), pp.109-160.\par\smallskip
\noindent [87]\ Eguchi, T., Hori, K. and Xiong, C., Quantum cohomology and Virasoro algebra, Physics Lett. B402(1997), pp. 71-80. hep-th/9703068.\par\smallskip
\noindent [88]\ Eguchi, T. and Xiong, C., Quantum cohomology at higher genus: topologi — cal recursion relations and Virasoro conditions. Adv. Theor. Math. Phys. 2(1997), pp. 219-229. hep-th/9801010.\par\smallskip
\noindent [89]\ Ekeland, I . , Convexity Methods in Hamiltonian Mechanics, Ergebnisse der Mathematik und ihrer Grenzgebiete(3) 19, Springer, Berlin, 1990,\par\smallskip
\noindent [90]\ Eliashberg, Y., "Cobordisme des solutions de relations differentielles" in South Rhone Seminar on Geometry, Travaux en Cours, Hermann, 1984.\par\smallskip
\noindent [91]\ Eliashberg, Y., A theorem on the structure of wave fronts and its applications in symplectic topology, Functional Anal. Appl. 21(1987), pp. 65-72.\par\smallskip
\noindent [92]\ Eliashberg, Y., "Filling by holomorphic disks and its applications" in Geometry of Low-Dimensional Manifolds 2 (Durham), Cambridge Univ. Press, 1990, pp. 45--68.\par\smallskip
\noindent [93]\ Eliashberg, Y., Contact 3-manifolds, twenty years since J. Martinet's work, Ann. Inst. Fourier 42(1992), pp.165-192.\par\smallskip
\noindent [94]\ Eliashberg, Y., "Invariants in contact topology" in International Congress of Math., 1998, Berlin. Documenta Mathematica, extra volume. http://www.mathematik.uni-bielefeld.de/ documenta.\par\smallskip
\noindent [95]\ Eliashberg, Y. and Gromov, M., Convex syrnplectic manifolds, Proc. Sympos. Pure Math. 52(1991), pp.135-162.\par\smallskip
\noindent [96]\ Eliashberg, Y. and Gromov, M., Lagrangian intersection and stable Morse the。ry, Boll. Un. Mat. Ital. B 11(1997), pp. 289-326.\par\smallskip
\noindent [97]\ Eliashberg, Y. and Thurston, W., Confoliations, University Lecture Series 13, Amer. Math. Soc., 1998.\par\smallskip
\noindent [98]\ Federer, H., Geometric Measure Theory, Die Grundlehren der mathematischen Wissenschaften in Einzeldarstellungen 153, Springer, Berlin, 1969.\par\smallskip
\noindent [99]\ Fedosov, B., A simple geometrical construction of deformation quantization, J . Di[!. Geom. 40(1994), pp. 213-238. [loo] Fedosov, B., Deformation Quantization and IndeのTheory, Mathematical Topics 9, Akadernie Verlag, 1996.\par\smallskip
\noindent [101]\ Floer, A., Proof of Arnold conjecture for surfaces and generalizations to certain 邸hler manifolds, Duke Math. J . 53(1986), pp. 1-32.\par\smallskip
\noindent [102]\ Floer, A., The unregularized gradient flow of the syrnplectic action, Comm. Pure Appl. Math. 41(1988), pp. 775-813.\par\smallskip
\noindent [103]\ Floer, A., Morse theory for Lagrangian intersections, J . Di[!. Geom. 28(1988), pp. 513-547.\par\smallskip
\noindent [104]\ Floer, A., An instanton invariant for 3-manifolds, Comm. Math. Phys. 118(1988), pp. 215-240.\par\smallskip
\noindent [105]\ Floer, A., Witten's complex and infinite dimensional Morse theory, J . Diff. Geom. 30(1989), pp. 207-221.\par\smallskip
\noindent [106]\ Floer, A., Symplectic fixed point and holomorphic spheres, Comm. Math. Phys. 120(1989), pp. 575-611.\par\smallskip
\noindent [107]\ Floer, A., Cup length estimate for Lagrangian intersections, Comm. Pure Appl. Math. 47(1989), pp. 335-356.\par\smallskip
\noindent [108]\ Floer, A. and Hofer, H., Coherent orientation problems for periodic orbits problem in symplectic geometry, Math. Z. 212(1993), pp.13-38.\par\smallskip
\noindent [109]\ Floer, A., Hofer, H. and Salamon, D., Transversality in elliptic Morse theory for the symplectic action, Duke Math. J . 80(1995), pp. 251-292.\par\smallskip
\noindent [110]\ Freed, D., Five Lectures on Super Symmetry, Amer. Math. Soc., 1999.\par\smallskip
\noindent [111]\ Friedman, R. and Morgan, J., Holomorphic principal bundles over elliptic curves. math/9811130. [119] Friedman, R., Morgan, J. and Witten, E., Vector bundles over elliptic fibrations, J . Alg. Geometry 8 (1999), pp. 279-401. alg-geom/9709029.\par\smallskip
\noindent [113]\ Friedman, R., Morgan, J . and Witten, E., Principal G-bundles over elliptic curves. Math. Res. Lett. 5 (1998), pp. 97-118. alg-geom/9707004.\par\smallskip
\noindent [114]\ Fukaya, K., "Morse homotopy, A00 category, and Floer homologies" in Gare Workshop on Geometry and Topology, Seoul National University, 1993.\par\smallskip
\noindent [115]\ Fukaya, K.,' 、Informal note on topology, geometry and topological field theory" in Geometry from the Pasific Rim, Herrick, A., Loo, B. and Wang, H.(eds.), Walter de Gruyter, 1994, pp. 99-116.\par\smallskip
\noindent [116]\ Fukaya, K., Morse homotopy and Chern-Simons peヰurbation theory, Comm. Math. Phys. 181(1996), pp. 37-90. [11 オ Fukaya, K., Floer homology of connected sum of homology 3-spheres, Topology 35(1996), pp. 89-136.\par\smallskip
\noindent [118]\ Fukaya, K., "Morse homotopy and its quantization" in Geometry and Topology, Kazez, W.(ed.), AMS/IP Studies in Advanced Math. 2-1, 1997, International Press, Hong Kong, pp. 409-440.\par\smallskip
\noindent [119]\ Fukaya, K., Floer homology for 3 manifold with boundary I , 1997, preprint, never to appear.\par\smallskip
\noindent [120]\ Fukaya, K., Floer homology of Lagrangian foliations and noncommutative mirror symmetry, preprint.\par\smallskip
\noindent [121]\ Fukaya, K., Mirror symmetry of Abelian variety and multi theta functions, J. Algebraic Geom. 11(2002), no. 3, pp. 393-512.\par\smallskip
\noindent [122]\ Fukaya, M., Oh, Y., Ohta, H. and Ono, K., Lagrangian intersection Floer theory: Anomaly and obstruction, preprint (2000).\par\smallskip
\noindent [123]\ Fukaya, K. and Oh, Y., Zero-loop open strings in the cotangent bundle and Morse homotopy, Asian. J. Math. 1(1997), pp. 96-180.\par\smallskip
\noindent [124]\ Fukaya, K. and Ono, K., Arnold conjecture and Gromov-Witten invariant, Topology 38(1999), pp, 933-1048.\par\smallskip
\noindent [125]\ Fukaya, F., Floer homology and mirror symmetry. I . Winter School on Mirror Symmetry, Vector Bundles and Lagrangian Submanifolds (Cambridge, MA, 1999), pp.15-43, AMS/IP Stud. Adv. Math., 23, Amer. Math. Soc., Providence, RI, 2001.\par\smallskip
\noindent [126]\ 深谷賢治，『ゲージ理論とトポロジ一』（シュプリンガー現代数学シリーズ），シュ プリンガー東京， 1996.\par\smallskip
\noindent [127]\ 深谷賢治，『解析力学と微分形式』（シリーズ「現代数学への入門」），岩波書店， 2004.\par\smallskip
\noindent [128]\ 深谷賢治，『これからの幾何学』， 日本評論社， 1998.\par\smallskip
\noindent [129]\ 深谷賢治編，『共形場の理論と幾何学』(Surveys in Geometry 予稿集）， 1992. http://www.math.nagoya-u.ac.jp/-naito/Surveys/List.html よりダウンロード可能\par\smallskip
\noindent [130]\ 深谷賢治編，『シンプレクティック幾何学』(Surveys in Geometry 予稿集）， 1995.\par\smallskip
\noindent [131]\ 深谷賢治編，『サイバーグ·ウィッテン理論』(Surveys in Geometry 予稿集）， 1997. http://www.ms.u-tokyo.ac.jp ; \textasciitilde{} furuta/index.html よりダウンロード可能\par\smallskip
\noindent [132]\ 深谷賢治編，『C' 環と幾何学』(Surveys in Geometry 予稿集）， 1998-\par\smallskip
\noindent [133]\ Fulton, W., Introduction to Torie Variety, Annals of Math. Studies 131, Princeton Univ. Press, 1993.\par\smallskip
\noindent [134]\ Fulton, W. and MacPherson, R., Compactification of configuration spaces, A叫 of Math. 139(1994), pp.183-225.\par\smallskip
\noindent [135]\ 古田幹雄，『指数定理』，岩波書店， 2008.\par\smallskip
\noindent [136]\ Gelfand, S. and Manin, I . , Method of Homological Algebra, Springer, Berlin, 1991.\par\smallskip
\noindent [137]\ Gelfand, I . , Kapranov, M. and Zelevinsky, A., Discriminants, Resultants, and Multidimensional Determinants, Mathematics: Theory \& Applications, Birkhauser, Basel, 1994.\par\smallskip
\noindent [138]\ Gerstenhaber, M., The cohomology structure of an associative ring, Ann. of Math. 78(1963), pp. 267-288.\par\smallskip
\noindent [139]\ Getzler, E., Virasoro conjecture for Gromov-Witten invariants, 1998. Algebraic geometry: Hirzebruch 70 (Warsaw, 1998), pp.147-176, Contemp. Math., 241, Amer. Math. Soc., Providence, RI, 1999. math/9903147.\par\smallskip
\noindent [140]\ Getzler, E. and Jones, J., A00 algebra and cyclic bar complex, Illinois J. Math. 34(1990), pp. 256-283.\par\smallskip
\noindent [141]\ Gilberg, D. and Trudinger, N., Elliptic Partial Differential Equations of Second Order, Grundlehren der mathematichen Wissenschaften 224, Springer, Berlin, 1970.\par\smallskip
\noindent [142]\ Ginzburg, V. and Kapranov, M., Koszul duality for operads, Duke Math. J. 76(1994), pp. 203-272.\par\smallskip
\noindent [143]\ Givental, A., "The nonlinear Maslov index" in Geometry of Low-Dimensional Manifolds 2 (Durham), Cambridge Univ. Press, 1990, pp. 35-43.\par\smallskip
\noindent [144]\ Givental, A., "A symplectic fixed point theorem for toric manifolds" in The Floer Memorial Volume, Hofer, H. et al.(eds.), Birkhauser, Basel, 1995, pp.445-481.\par\smallskip
\noindent [145]\ Givental, A., Homological geometry, I : projective hypersurfaces, Selecta. Math. 1(1995), pp. 325-345.\par\smallskip
\noindent [146]\ Givental, A., Equivariant Gromov-Witten invariants, Internat. Math. Res. Notices 13(1996), pp. 613-663. [141] Givental, A., "Elliptic Gromov-Witten invariants and the generalized mirror conjecture" in Integral Systems and Algebraic Geometry, World Scientific, 1998, pp. 107-155.\par\smallskip
\noindent [148]\ 牛腸徹， "Connection 付きのhermitian line bundle をめぐって", 『共形場の理論 と幾何学』(Surveys in Geometry 予稿集深谷賢治編）所収， 1992-\par\smallskip
\noindent [149]\ 牛腸徹， "SUPERSYMMETRY 入門", 『サイバーグ·ウィッテン理論』(Surveys in Geometry 予稿集，深谷賢治編）所収， 1997.\par\smallskip
\noindent [150]\ Goldman, W., The symplectic nature of fundamental groups of surfaces, Adv. in Math. 54(1984), pp. 200-225.\par\smallskip
\noindent [151]\ Goldman, W., Invariant functions on Lie groups and Hamiltonian flows of surface group representations, Invent. Math. 85(1986), pp. 263-302.\par\smallskip
\noindent [152]\ Goldman, W. and Millson, J., The deformation theory of representations of fundamental groups of compact Kahler manifolds, Publ. Math. !HES 67(1988), pp. 43-96. 瞬! Goldman, W. and Millson, J., The homotopy invariance of the Kuranishi space, fllinois J. Math. 34(1990), pp. 337-367. [1541 Gompf, R., A new construction of symplectic manifolds, Ann. of Math. 142(1995), pp. 527-595. 瞑I Green, M., Schwarz, J. and Witten, E., Super String Theory I,II, Cambridge Monographs on Mathematical Physics, Cambridge Univ. Press, 1987. [156) Griffiths, P. and Harris, J., Principle of Algebmic Geometry, John Wiley, 1978. [15竹 Griffiths, P. and Morgan, J., Rational Homotopy Theory and Differential Forms, Progress\_in Math. 16, Birkhauser, Basel, 1981. 且58] Gromov, M., Filling Riemannian manifolds, J. Diff. Geom. 18(1983), pp.1-147. 飼9] Gromov, M., Pseudo holomorphic curves in symplectic manifolds, Invent. Math. 82(1985), pp. 307ー347. (\par\smallskip
\noindent [160]\ Gromov, M., Partial Differential Relations, Ergebnisse der Mathematik und ihrer Grenzgebiete 3-9, Springer, Berlin, 1986.\par\smallskip
\noindent [161]\ Gromov, M., "Soft and hard symplectic geometry" in International Congress of Mathematics, 1986, Berkeley. (\par\smallskip
\noindent [162]\ Gross, M. and Wilson, P., Mirror symmetry via 3-tori for a class of Calabi-Yau threefolds, Math. Ann. 309(1997), pp. 505-531. 邸3] \_ Gross, M., "Special Lagrangian fibrations, I , topology" in Integmble Systems and Algebmic Geometry, World Scientific, River Edge, 1997, pp.156-193. [1641 Gross, M., Special Lagrangian fibrations II: Geometry. Winter School on Mirror Symmetry, Vector Bundles and Lagrangian Submanifolds (Cambridge, MA, 1999), pp. 95-150, AMS/IP Stud. Adv. Math., 23, Amer. Math. Soc., Providence, RI, 2001. math/9809072. [165) Guillemin, V. and Sternberg, S., Geometric Asymptotics, Mathematical Surveys 14, Amer. Math. Soc., 1977.\par\smallskip
\noindent [166]\ Guillemin, V. and Sternberg, S., Convexity properties of the moment map, Invent. Math. 97(1982), pp. 485-522.\par\smallskip
\noindent [167]\ Guillemin, V. and Sternberg, S., Symplectic Techniques in Physics, Cambridge Univ. Press, 1984. [1681 Guillemin, V. and Sternberg, S., On the Kostant multiplicity formula, J. Geom. and Phys. 5(1989), pp. 721-750. [169) Guillemin, V. and Sternberg, S., Supersymmetry and Equivariant de Rham The。ry, Springer, Berlin, 1999.\par\smallskip
\noindent [170]\ Giittsche, L. and Zagier, D., Jacobi forms and the structure of Donaldson invariants for 4-manifolds with 如= 1, Selecta Math. 4(1998), pp. 69-115. (\par\smallskip
\noindent [171]\ 原田耕一郎，『モンスター』，岩波書店， 1999.\par\smallskip
\noindent [172]\ Hardt, R. and Simon, L., Seminar on Geometric Measure Theory, DMV Seminar, 7, Birkhauser, Basel, 1986.\par\smallskip
\noindent [173]\ Hartshorne, R., Residues and Duality, Lecture Note in Math. 20, Springer, Berlin, 1966. [1741 Harvey, R. and Lawson, B., Calibrated geometries, Acta Math. 148(1982), pp.47-157.\par\smallskip
\noindent [175]\ 服部晶夫，『位相幾何学』（基礎数学選書），岩波書店， 1991-\par\smallskip
\noindent [176]\ 服部晶夫，『多様体のトポロジ一』，岩波書店， 2003- [177'1 Herald, C., Legendrian cobordism and Chern-Simons theory on 3-manifolds with boundary, Comm. Anal. Geom. 3(1994), pp. 337--413.\par\smallskip
\noindent [178]\ Hiilmander, L., Hypoelliptic second order differential equations, Acta Math. 119(1967), pp.147-171.\par\smallskip
\noindent [179]\ Hirsh, M., Immersions of manifolds, 加ns. Amer. Math. Soc. 93(1961), pp. 242-276.\par\smallskip
\noindent [180]\ Hitchin, N., Karlhede, A., Lindstrom, U. and Rocek, M., Hyper-Kahler metrics and supersymmetry, Comm. Math. Phys. 108(1987), pp. 535-589.\par\smallskip
\noindent [181]\ Hofer, H., Pseudoholomorphic curves in symplectisations with application to the Weinstein conjecture in dimension three, Invent. Math. 114(1993), pp. 515-563. [18乱 Hofer, H. and Zehnder, E., Symplectic Invariants and Hamiltonian Dynamics, Birkhふuser Advanced Texts, Birkhauser, Basel, 1994.\par\smallskip
\noindent [183]\ Hofer, H. and Salamon, D., "Floer homology and Novikov ring" in The Floer Memorial Volume, Hofer, H. et al.(eds.), Birkh細user, Basel, 1995, pp. 483-524.\par\smallskip
\noindent [184]\ Hummel, C., Gromov's Compactness Theorem for Pseudo-holomorphic Curves, Progress in Math. 151, Birkhauser, Basel, 1997.\par\smallskip
\noindent [185]\ Hutchings, M., Reidemeister torsion in generalized Morse theory, Thesis, De— partment of Mathematics, Harvard Univ., 1998.\par\smallskip
\noindent [186]\ Hutchings, M. and Lee, Y., Circle valued Morse theory, Reidemeister torsion, and Seiberg-Witten invariants of three manifolds, Topology 38(1999), pp. 861-888. [1811 Itagaki, Y., Convergence of the normalized solution of the Maurer-Cartan equa-- tion in the Barannikov-Kontsevich construction, Master Thesis, Kyoto Univ, 1999.\par\smallskip
\noindent [188]\ 泉屋周ー•石JI I剛郎，『応用特異点論』，共立出版， 1998,\par\smallskip
\noindent [189]\ Jeffrey, L. and Kirwan, F., Localization for nonabelian group actions, Topology 34(1995), pp. 291-327.\par\smallskip
\noindent [190]\ Jeffrey, L. and Kirwan, F., Intersection theory on moduli spaces of holomorphic bundles of arbitrary rank on a Riemann surface, Ann. of Math. 148(1998), pp. 109.,\_ 196.\par\smallskip
\noindent [191]\ Jeffrey, L. and Weitsman, J., Half density quantization of the moduli space of flat connections and Witten's semiclassical manifold invariants, Topology 32(1993), pp. 509-529.\par\smallskip
\noindent [192]\ Jeffrey, L. and Weitsman, J., Bohr-Sommerfeld orbits in the moduli space of flat connections and the Verlinde dimension formula, Comm. Math. Phys. 150(1993), pp. 593--630.\par\smallskip
\noindent [193]\ Joyce, D., Compact Riemannian 7-manifolds with holonomy G2, I , II, J. Diff. Geom. 43(1996), pp. 291-328, 359-375.\par\smallskip
\noindent [194]\ Joyce, D., Compact 8-manifolds with holonomy Spin(7), Invent. Math. 123(1996), pp. 507-552.\par\smallskip
\noindent [195]\ Kapranov, M., "Operads and algebraic geometry" in International Congress of Mathematics, 1998, Berlin.\par\smallskip
\noindent [196]\ Karasev, M. and Maslov, V., Nonlinear Poisson Brackets, Geometry and Quantization, Translations of Mathematical Monographs 119, Amer. Math. Soc., 1993.\par\smallskip
\noindent [197]\ Keel, S., Intersection theory of the moduli space of stable n-pointed curves of genus zero, Trans. Amer. Math. Soc. 330(1992), pp. 545-574.\par\smallskip
\noindent [198]\ Kimura, T., Stasheff, J., Voronov, J. and Alexander, A., On operad structures of moduli spaces and string theory, Comm. Math. Phys. 171(1995), pp. 1-25.\par\smallskip
\noindent [199]\ Kirillov, "Geometric quantization" in Encyclopaedia of Mathematical Sciences, Dynamical System IV, Arnold, V. and Novikov, S.(eds.), Springer, Berlin, 1980, pp. 254-272.\par\smallskip
\noindent [200]\ Kirwan, F., Cohomology of Quotients in Symplectic and Algebraic Geometry, Mathematical Notes 31, Princeton Univ. Press, 1984.\par\smallskip
\noindent [201]\ Kirwan, F., The cohomology rings of moduli spaces of vector bundles over Riemann surfaces, J. Amer. Math. Soc. 5(1992), pp. 853-906.\par\smallskip
\noindent [202]\ Knudsen, F., Projectivity of the moduli space of stable curves, II: the stacks M 。,n, Math. Scand. 52(1983), pp.1225-1265.\par\smallskip
\noindent [203]\ Kobayashi, R. and Todorov, A., Polarized period map for generalized K3 surfaces and the moduli of Einstein metric, Tohoku Math. J. 39(1987), pp. 341-363.\par\smallskip
\noindent [204]\ Kobayashi, S. and Nomizu, K., Foundation of Differential Geometry I,II, Inter-science Tracts in Pure and Applied Mathematics 15, Interscience, 1963, 1969.\par\smallskip
\noindent [205]\ 小林昭七，『接続の微分幾何とゲージ理論』，裳華房， 1989.\par\smallskip
\noindent [206]\ 小林昭七，『複素幾何』，岩波書店， 2005.\par\smallskip
\noindent [207]\ Kodaira, K. and Spencer, D., On deformation of complex analytic structures I , II, Ann. of Math. 67(1958), pp. 328-466.\par\smallskip
\noindent [208]\ Kodaira, K., Nirenberg, L. and Spencer, D., On the existence of deformations of complex analytic structures, Ann. of Math. 68(1958), pp. 450-459.\par\smallskip
\noindent [209]\ 小平邦彦，複素多様体と複素構造の変形1, 2, 東大セミナリーノート 19, 31, 東 京大学数学教室， 1974.\par\smallskip
\noindent [210]\ 小平邦彦，『複素多様体論』，岩波書店， 1992-\par\smallskip
\noindent [211]\ 河野俊丈，『曲面の幾何構造とモジュライ』， 日本評論社， 1997.\par\smallskip
\noindent [212]\ 河野俊丈，『場の理論とトポロジ一』，岩波書店， 2008-\par\smallskip
\noindent [213]\ Kolmogorov, A., On the representation of continuous functions of several variables as super position of functions of smaller number of variables, Soviet Math. Dokl. 108(1956), pp. 179-182. (Selected works of Kolmogorov, Vol. 1 , pp. 378-382)\par\smallskip
\noindent [214]\ Konno, H., On the natural line bundle on the moduli space of stable parabolic bundles, Comm. Math. Phys. 155(1993), pp. 311-324.\par\smallskip
\noindent [215]\ Kontsevich, M., "A00-algebras in mirror symmetry" in Arbeitstagung, Max Plank Institute of Mathematics, 1992.\par\smallskip
\noindent [216]\ Kontsevich, M., "Formal (non)commutative symplectic geometry" in The Gelfand Mathematical Seminars 1990-1992, Birkhiiuser, Zi.irich, 1993, pp.173-187.\par\smallskip
\noindent [217]\ Kontsevitch, M., "Feynman diagram and low dimensional topology" in Proceeding of the First European Congress of Mathematics, Birkhiiuser, Boston, 1994, pp. 97-122.\par\smallskip
\noindent [218]\ Kontsevich, M., "Enumeration of rational curve by torus action" in Moduli Space of Surface, Dijkgraaf, H., Faber, C. and v. d. Geer, G.(eds.), Birkhiiuser, ;3oston, 1995, pp. 335-368.\par\smallskip
\noindent [219]\ Kontsevich, M., "Homological algebra of mirror symmetry" in International Congress of Mathematics, Birkhiiuser, Zi.irich, 1995.\par\smallskip
\noindent [220]\ Kontsevich, M., "Formality conjecture" in Deformation Theory and Symplec— tic Geometry, Steinheimer, D., Rawnsley, J . and Gutt, S.(eds.), Kluwer Academic Publishers, 1997, pp.139-156.\par\smallskip
\noindent [221]\ Kontsevich, M., Deformation quantization of Poisson manifolds I , Lett. Math. Phys. 66(2003), no. 3, 157-216. q-alg/9709040.\par\smallskip
\noindent [222]\ Kontsevich, M. and Manin, Y., Gromov-Witten classes, quantum cohomology, and enumerative geometry, Comm. Math. Phys. 164(1994), pp. 525-562. hepth/9102147.\par\smallskip
\noindent [223]\ Kostant, B., Quantization and unitary representations, Lecture Notes in Math. 170, Springer, Berlin, 1970, pp. 87-208. [2241 Kotschick, D., The Seiberg-Witten invariants of symplectic four-manifolds (after C.H. Taubes), Seminaire Bourbaki 1995/96, Asterisque 241, 1997, pp.195-220.\par\smallskip
\noindent [225]\ Kronheimer, P., The construction of ALE spaces as hyper Kahler quotients, J. Diff. Geom. 29(1989), pp. 665-683.\par\smallskip
\noindent [226]\ Kronheimer, P. and Mrowka, T., Monopoles and contact structures, Invent. Math. 130(1997), pp. 209-255. 加1J Kronheimer, P. and Nakajima, H., Yang-Mills instantons on ALE gravitational instantons, Math. Ann. 288(1990), pp. 263-307.\par\smallskip
\noindent [228]\ Kuksin, F., Infinite-dimensional symplectic capacities and a squeezing theorem for Hamiltonian PDEs, Comm. Math. Phys. 167(1995), pp. 531-552. 位29] Kuranishi, M., On the locally complete families of complex a叫ytic structures, Ann. of Math. 75(1962), pp. 536-577. 位30] Lalonde, F. and McDuff, D., "J-curves and the classification of rational and ruled symplectic 4-manifolds" in Contact and Symplectic Geometry, Thomas, C.(ed.), Cambridge Univ. Press, 1996. 位31] Lalonde, F., McDuff, D. and Polterovich, L., "〇n the flux conjectures" in Geometry, Topology, and Dynamics, CRM Proc. Lecture Notes 15, Amer. Math. Soc., 1998, pp. 69-85. 似32] Lees, J., On the classification of Lagrange immersion, Duke Math. J. 43(1976), pp. 217-224. 加! Lian, B., Liu, K. and Yau, S. T., Mirror principle I , Asian J. Math. 1(1998), pp. 729-763. [234! Lie, J. and Tian, G., Virtual moduli cycles and Gromov Witten invariants of algebraic varieties, J. Amer. Math. Soc. 6(1997), pp. 269-305. 邸J Lie, J. and Tian, G., "Virtual moduli cycles and Gromov Witten invariants of general symplectic manifolds" in Topics in Symplectic 4-Manifolds, International Press, Cambridge, 1998, pp. 47-83.\par\smallskip
\noindent [236]\ Liu, G. and Tian, G., Floer homology and Arnold conjecture, J . Diff. Geom. 49(1998), pp.1-74. 松37] Liu, X. and Tian, G., Virasoro Constraints for Quantum cohomology. J. Differential Geom. 50(1998), no. 3, 537-590. math/9806028.\par\smallskip
\noindent [238]\ Loday, J., Stasheff, J., and Voronov, A.(eds.), Operads: proceedings of renais — sance conferences, Contemporary Math. 102 (1997). ね39] Lyusternik, L. and Snirel'man, L., Methodes Topologiques dans les Problもmes Variationnels, Actualites Sci. Indust. 118, Hermann, 1934.\par\smallskip
\noindent [240]\ MacLane, S., Homology, Die Grundlehren der mathematischen Wissenchaften in Einzeldarstellungen 114, Springer, Berlin, 1967. 松41] Manin, Y., "Generating functions in algebraic geometry and sums over trees" in Moduli Space of Surface, Dijkgraaf, H., Faber, C. and v. d. Geer, G.(eds.), Birkhauser, Boston, 1995, pp. 401-418.\par\smallskip
\noindent [242]\ Manin, Y., Frobenius Manifolds, Quantum Cohomology, and Moduli Spaces, Colloquim Pub!. 47, Amer. Math. Soc., 1999. 似3] Manin, Y, Three constructions of Frobenius manifolds: a comparative study, Surveys in differential geometry, pp. 497-554, Surv. Differ. Geom., VII math/9801006.\par\smallskip
\noindent [244]\ Marsden, J., "Park City Lectures on Mechanics, Dynamics, and Symmetry" in Symplectic Geometry and Topology, Eliashberg, Y. and Traynor, L. (eds.), IAS/Park City Math. Series 7, Amer. Math. Soc. 1999, pp. 335-430.\par\smallskip
\noindent [245]\ Marsden, J. and Weinstein, A., Reduction of symplectic manifolds with symme— try, Reports on Mathematical Physics 5(1974), pp.121-130.\par\smallskip
\noindent [246]\ マスロフ，『摂動論と漸近的方法』，大内忠·金子晃·村田実訳，岩波書店， 1965.\par\smallskip
\noindent [247]\ Mathai, V. and Quillen, D., Thom classes, super connections, and equivariant differential forms, Topology 25(1986), pp. 85-110. 似48] 松本幸夫，『Morse 理論の基礎』，岩波書店， 2005. 似49] May, J., The geometry of iterated loop spaces, Lecture Note in Math. 271, Springer, Berlin, 1972.\par\smallskip
\noindent [250]\ McDuff, D., Example of simply connected symplectic non Kahlerian manifolds, J. Diff. Geom. 77(1984), pp. 267-277. 肱51] McDuff, D., Example of symplectic structures, Invent. Math. 89(1987), pp. 13-36.\par\smallskip
\noindent [252]\ McDuff, D., Elliptic methods in symplectic geometry, Bull. Amer. Math. Soc. 23(1990), pp. 311-358.\par\smallskip
\noindent [253]\ McDuff, D. and Salamon, D., J-holomorphic Curves and Quantum Cohomology, University Lecture Series 6, Amer. Math. Soc., 1994.\par\smallskip
\noindent [254]\ McDuff, D. and Salamon, D., Introduction to Symplectic Topology, Oxford Mathematical Monographs, Oxford Univ. Press, 1995.\par\smallskip
\noindent [255]\ McLean, R., Deformations of calibrated submanifolds, Comm. Anal. Geom. 6(1998), pp. 705-747.\par\smallskip
\noindent [256]\ Milnor, J., On the cobordism ring 几and its complex analogue, Amer. J. Math. 82(1960), pp. 505-521.\par\smallskip
\noindent [257]\ Milnor, J., Morse Theory, Annals of Math. Studies 51, Princeton Univ. Press, 1963. (邦訳）『モース理論』，志賀浩二訳，吉岡書店， 1983-\par\smallskip
\noindent [258]\ Milnor, J., Lectures on the h-Cobordism Theorem, Mathematical Notes, Princeton Univ. Press, 1965.\par\smallskip
\noindent [259]\ Milnor, J., Whitehead torsion, Bull. Amer. Math. Soc. 72(1966), pp. 358-426.\par\smallskip
\noindent [260]\ Minahan, J., Nemeschansky, D., Vafa, C. and Warner, N., E-strings and N = 4 topological Yang-Mills theories, Nucl. Phys. B517(1998), pp. 537-556. hepth/9802168.\par\smallskip
\noindent [261]\ Moore, G. and Witten, E., Integration over the u-plane in Donaldson theory, Adv. Theor. Math. Phys. 1(1997), pp. 298-387. hep-th/9709193.\par\smallskip
\noindent [262]\ Morgan, J., Mrowka, T. and Ruberman, D., The2 Moduli Space and a Vanishing Theorem for Donaldson Polynomial Invariants, Monographs in Geometry and Topology II, International Press, 1994.\par\smallskip
\noindent [263]\ Kollar, J. ·森重文，『双有理幾何学』，岩波書店， 2008.\par\smallskip
\noindent [264]\ 森田茂之，『特性類と幾何学』，岩波書店， 2008.\par\smallskip
\noindent [265]\ Moser, J., On the volume elements on manifolds, Trans. of Amer. Math. Soc. 120(1965), pp. 280-296.\par\smallskip
\noindent [266]\ Mukai, S., Duality between D(X) and D(X) with its application to Picard sheaves, Nagoya Math. J. 81(1981), pp. 334-353.\par\smallskip
\noindent [267]\ Mukai, S., Abelian variety and spin representation, preprint, 1998. Translation of an article written in Japanese in 1994.\par\smallskip
\noindent [268]\ 向井茂，『モジュライ理論1, 2 』，岩波書店， 2008.\par\smallskip
\noindent [269]\ Mumford, D., Geometric Invariant Theory, Springer, Berlin, 1965.\par\smallskip
\noindent [270]\ Mumford, D., Abelian Variety, Oxford Univ. Press, 1974.\par\smallskip
\noindent [271]\ Mumford, D., Stability of projective varieties, Enseign. Math. 24(1977), pp. 39-110.\par\smallskip
\noindent [272]\ Mumford, D., "Toward an enumerative geometry of the moduli space of curves" in Arithmetric and Geometry, Artin, M. and Tate, J.(eds.), Birkhauser, Boston, 1983, pp. 271ー326.\par\smallskip
\noindent [273]\ 村井信行，『拘束系の力学』， 日本評論社， 1998-\par\smallskip
\noindent [274]\ 長野正，『大域変分法』（現代の数学17), 共立出版， 1971.\par\smallskip
\noindent [275]\ Nakajima, H., Removable singularities for Yang-Mills connections in higher dimensions, J. Fae. Sci. Univ. Tokyo Sect. IA Math. 34(1987), pp. 299-307.\par\smallskip
\noindent [276]\ Nakajima, H., Compactness of the moduli space of Yang-Mills connections in higher dimensions, J. Math. Soc. Japan, 40(1988), pp. 383-392.\par\smallskip
\noindent [277]\ Nakajima, H., Hausdorff convergence of Einstein 4-manifolds, J. Fae. Sci. U血 Tokyo Sect. IA Math. 35(1988), pp. 411-424.\par\smallskip
\noindent [278]\ Nakajima, H., Instantons on ALE spaces, quiver varieties, and Kac-Moody algebras, Duke Math. J. 2(1994), pp. 61-74.\par\smallskip
\noindent [279]\ Nakajima, H., Heisenberg algebra and Hilbert schemes of points on projective surfaces, Ann. of Math. 145(1997), pp. 379-388.\par\smallskip
\noindent [280]\ Nakajima, H., Lectures on Hilbert scheme of points on surface, University Lee— ture Series, 18, Amer. Math. Soc., Providence, RI, 1999.\par\smallskip
\noindent [281]\ 中島啓『非線形問題と複素幾何学』，岩波書店， 2008.\par\smallskip
\noindent [282]\ 中岡稔『位相幾何学 ホモロジー論—』（現代の数学15), 共立出版， 1960.\par\smallskip
\noindent [283]\ Narashimhan, M. and Seshadri, C., Stable and unitary vector bundles on compact Riemannian surfaces, Ann. of Math. 82(1965), pp. 540-567.\par\smallskip
\noindent [284]\ Newlander, R. and Nirenberg, L., Complex analytic coordinates in almost complex manifolds, Ann. of Math. 65(1957), pp. 391-404.\par\smallskip
\noindent [285]\ 西川青季『幾何学的変分問題』，岩波書店， 2006.\par\smallskip
\noindent [286]\ Novikov, S., Solitons and Geometry, Academia Nazionale dei Lincei, Scuola Normale Superiore, Press Syndicate of the University of Cambridge, 1992. [28竹 Novikov, S., Multivalued functions and functionals - an analogue of the Morse theory, Soviet Math. Dokl. 24(1981), pp. 222-225.\par\smallskip
\noindent [288]\ 小田忠雄，『凸体と代数幾何学』（紀伊國屋数学叢書24), 紀伊國屋書店， 1988.\par\smallskip
\noindent [289]\ Oh, Y., Removal of boundary singularities of pseudo-holomorphic curves with Lagrangian boundary conditions, Comm. Pure Appl. Math. 45(1992), pp.131-139.\par\smallskip
\noindent [290]\ Oh, Y., Floer cohomology of Lagrangian intersections and pseudo-holomorphic disks I , II, Comm. Pure and Appl. Math. 46(1993), pp. 949-994, 995-1012.\par\smallskip
\noindent [291]\ Oh, Y., Floer cohomology, spectral sequences, and the Maslov class of Lagrangian embeddings, Internat. Math. Res. Notices 7(1996), pp. 305ー346.\par\smallskip
\noindent [292]\ 岡本久·中村周，『関数解析』，岩波書店， 2006.\par\smallskip
\noindent [293]\ Ono, K., On the Arnold conjecture for weakly monotone symplectic manifolds, 加ent. Math. 119(1995), pp. 519-537.\par\smallskip
\noindent [294]\ Omori, H., Maeda, Y. and Yoshioka, A., Weyl manifolds and deformation quan-tization, Adv. Math. 85(1991), pp. 224-255.\par\smallskip
\noindent [295]\ 大森秀樹『無限次元Lie 群論』（紀伊國屋数学叢書15), 紀伊國屋書店， 1978-\par\smallskip
\noindent [296]\ 大森秀樹，『一般力学系と場の幾何学』，裳華房， 1991. [29オ 大沢健夫，『多変数複素解析』，岩波書店， 2008.\par\smallskip
\noindent [298]\ Pansu, P., "Compactness" in Holomorphic Curves in Symplectic Geometry, Audin, M. and Lafontaine, J.(eds.), Birkhauser, Basel, 1994, pp. 233—250.\par\smallskip
\noindent [299]\ Parker, T. and Wolfson, J., Pseudo holomorphic maps and bubble trees, J. Geom. Analysis 3(1993), pp. 63-98.\par\smallskip
\noindent [300]\ Pazhitnov, Rationality of boundary operators in the Novikov complex in general position, St. Petersburg Math. J. 9(1998), pp. 969-1006.\par\smallskip
\noindent [301]\ Pazhitnov, Incidence coefficients in the Novikov complex for Morse forms: ra— tionality and exponential growth properties. dg-ga/9604004.\par\smallskip
\noindent [302]\ Piunikihin, S., Schwartz, M. and Salamon, D., "Symplectic Floer-Donaldson theory and quantum homology" in Contact and Symplectic Geome四，Thomas, C.(ed.), Cambridge Univ. Press, 1996, pp. 171-200.\par\smallskip
\noindent [303]\ Polchinski, J., "TASI lectures on D-branes" in Fields, String, and Duality, Efthimiou, C. and Green, B.(eds.), World Scientific, Hong-Kong, 1997. hepth/9611050.\par\smallskip
\noindent [304]\ Polchinski, J., String Theory I , II, Cambridge Univ. Press, 1998.\par\smallskip
\noindent [305]\ Polishchuk, A. and Zaslow, E., Categorical mirror symmetry: the elliptic curve, Adv. Theor. Math. Phys., 2(1998), pp. 443-470. math/9801119.\par\smallskip
\noindent [306]\ Polishchuk, A., Massey and Fukaya product on Elliptic curves. math.AG/980301.\par\smallskip
\noindent [307]\ Polterovich, L., The surgery of Lagrange submanifolds, Geometric Analysis and Functional Analysis, 1(1991), pp.198-210.\par\smallskip
\noindent [308]\ Polterovich, L., Monotone Lagrange submanifolds of linear spaces and the Maslov class in cotangent bundles, Math. Z. 207(1991), pp. 217-222.\par\smallskip
\noindent [309]\ Pozniak, M., Floer homology, Novikov rings, and clean intersections, Thesis, Mathematics Institute, University of Warwick, Coventry, 1994.\par\smallskip
\noindent [310]\ Quillen, D., Rational homotopy theory, A皿 of Math. 90(1969), pp. 205-295.\par\smallskip
\noindent [311]\ Rabinowitz, P., Periodic solutions of Hamiltonian systems on a prescribed energy surface, J . Diff. Equations 33(1978), pp. 336-352.\par\smallskip
\noindent [312]\ Ran, Z., "Thickening Calabi-Yau moduli spaces" in M訂or Symmetry II, Greene, B. and Yau, S. T.(eds.), International Press, Hong-Kong, 1997, pp. 393-400.\par\smallskip
\noindent [313]\ Rieffel, M., C'algebras associated with irrational rotations, Pacific J . Math. 93(1981), pp. 415-429.\par\smallskip
\noindent [314]\ Rieffel, M., Noncommutative tori\textasciitilde{}a case study of noncommutative differential geometry, Contemporary Math. 105(1990), pp. 191-211.\par\smallskip
\noindent [315]\ Ruan, Y., Symplectic topology and extremal rays, Geometric and Functional Analysis 3(1993), pp. 395-430.\par\smallskip
\noindent [316]\ Ruan, Y., Topological sigma model and Donaldson-type invariants in Gromov theory, Duke Math. J . 83(1996), pp. 461-500.\par\smallskip
\noindent [317]\ Ruan, Y., Virtual neighborhood and pseudoholomorphic curve, Proceedings of 6th Gokova Geometry-Topology Conference. Turkish J. Math. 23(1999), no.1, pp. 161-231.\par\smallskip
\noindent [318]\ Ruan, Y. and Tian, G., Bott-type symplectic Floer cohomology and its multiplicative structure, Mathematical Research Letters 2(1995), pp. 203-219.\par\smallskip
\noindent [319]\ Ruan, Y. and Tian, G., A mathematical theory of quantum cohomology, J . Diff. Geom. 42(1995), pp. 259-367.\par\smallskip
\noindent [320]\ Sacks, J . and Uhlenbeck, K., The existence of minimal immersions of 2-spheres, A加. of Math. 113(1981), pp.1-24.\par\smallskip
\noindent [321]\ Saito, K., Period mapping associated to a primitive form, Publ. RIMS 19(1983), pp.1231-1264. [32乱 Saito, K., Duality of regular system of weights, Asian J . Math. 2(1998), no. 4, pp. 983-1047.\par\smallskip
\noindent [323]\ Satake, I., On a generalization of the notion of manifold, Proc. Nat. Acad. Sci. USA, 42(1956), pp. 359--363.\par\smallskip
\noindent [324]\ Schechtman, V., Remarks on formal deformations and Batalin-Vilkovsky algebras. math/9802006.\par\smallskip
\noindent [325]\ Schlessinger, M., Functors of Artin rings, Trans. Amer. Math. Soc. 130(1968), pp. 208-222. 似26] Shlessinger, M. and Stasheff, J., The Lie algebra structure on tangent cohomology and deformation theory, J. Pure Appl. Algebra 89(1993), pp. 231-235.\par\smallskip
\noindent [327]\ Shoen, R. and Wolfson, J., Minimizing volume among Lagrangian submanifolds, Proc. Sympos. Pure Math. 65(1998), pp.181-199. 肱28] Schwarz, M., Morse Homology, Progress in Math. 111, Birkhauser, Basel, 1993.\par\smallskip
\noindent [329]\ Seiberg, N. and Witten, E., Electric-magnetic duality, monopole condensa — tion, and confinement in N = 2 supersymmetric Yang-Mills theory, Nucl. Phys. B303(1994), p. 19. hep-th/9407087. 侶30] Seidel, P., 1r1 of symplectic automorphism groups and invertibles in quantum homology rings, Geometric and Functional Analysis 6(1997), pp.1046-1095.\par\smallskip
\noindent [331]\ Shaneson, J., "Characteristic classes, lattice points, and Euler-MacLaurin formulae" in International Congress of Mathematics, Birkhauser, Ziirich, 1994. [33乱 重川一郎，『確率解析』，岩波書店， 2008,\par\smallskip
\noindent [333]\ Siebert, B., Gromov-Witten invariants for general symplectic manifolds. alggeom/9611021.\par\smallskip
\noindent [334]\ Sikorav, J., Points fixes d'une application symplectique homologue l'identite, J . Diff. Geom. 22(1985), pp. 49-79.\par\smallskip
\noindent [335]\ Sikorav, J., Homologie de Novikov associe a une classe de cohomologie reelle de degre un, Thesis, Paris: Orsay, 1987. 似36] Silva., A.C. and Weinstein A., Geometric Models for Noncommutative Algebras, Berkeley Math., Lecture Notes 10, Amer. Math. Soc., 1999.\par\smallskip
\noindent [337]\ Silva, V. D., Products on symplectic Floer homology, Thesis, Oxford University, 1997. 1'\par\smallskip
\noindent [338]\ Simon, L., Asymptotics for a class of non-linear evolution equations with application to geometric problems, Ann. of Math. 118(1983), pp. 525-571.\par\smallskip
\noindent [339]\ Simon, L., Lectures on geometric measure theory, Proc. Centre Math. Anal. Austral. Nat. Univ. 3(1983). 1'\par\smallskip
\noindent [340]\ Simpson, C., Algebraic aspects of higher nonabelian Hodge theory. math/9902067.\par\smallskip
\noindent [341]\ Siu, Y., Lectures on Hermitian-Einstein metrics for stable bundles and K辿ler­ Einstein metrics, Birkhauser, Diisseldorf, 1986. [3421 Smale, S., An infinite dimensional version of Sard's theorem, Amer. J . Math. 87(1968), pp. 861-866.\par\smallskip
\noindent [343]\ Sniatycki, J., 0 n cohomology group appearmg in geometnc quant1zat10n, Lee-ture Note in Math. 570, Springer, 1977, pp. 46-66.\par\smallskip
\noindent [344]\ 紐iatycki, J., Geomet加Quantization and Quantum Mechanics, Applied Mathematical Sciences 30, Springer, Berlin, 1980.\par\smallskip
\noindent [345]\ Souriau, J., Quantification geometrique, Comm. Math. Phys. 1(1966), pp. 374-398.\par\smallskip
\noindent [346]\ Spanier, E., Algebraic Topology, McGraw Hill, 1966.\par\smallskip
\noindent [347]\ St邸heff, J., Homotopy 邸sociativity of H-Spaces I , II. Trans. Amer. Math. Soc. 108(1996), pp. 275-292, 293--312.\par\smallskip
\noindent [348]\ St邸heff, J., "Deformation theory and Batalin-Virkovsky Master equation" in Deformation Theory and Symplectic Geometry, Steinheimer, D., Rawnsley, J. and Gutt, S.(eds.), Kluwer Academic Publishers, 1997, pp.271-284.\par\smallskip
\noindent [349]\ Steenrod, N. and Epstein, D., Cohomology Operations, Ann. of Math. Studies 50, Princeton Univ. Press, 1962.\par\smallskip
\noindent [350]\ Strominger, A., Yau, S. T. and Z邸low, E., Mirror syillilletry i s T-duality, Nucl. Phys. B476(1996), pp. 243-259. hep-th/9606040.\par\smallskip
\noindent [351]\ Sullivan, D., Infinitesimal calculations in topology, Publ. Math. IHES 74(1978), pp. 269-331.\par\smallskip
\noindent [352]\ Swan, R., Vector bundles and projective modules, Trans. Amer. Math. Soc. 105(1962), pp. 264-277.\par\smallskip
\noindent [353]\ Takakura, T., Degeneration of Riemann surfaces and intermediate polarization of the moduli space of flat connections, Invent. Math. 123(1996), pp. 431-452.\par\smallskip
\noindent [354]\ 高倉樹，‘‘幾何学的量子化の理論概観", 『シンプレクティック幾何学』(Surveys in Geometry 予稿集，深谷賢治編）所収， 1995,\par\smallskip
\noindent [355]\ 田村一郎，『葉層のトボロジ一』（数学選書），岩波書店， 1976,\par\smallskip
\noindent [356]\ 田村一郎，『微分位相幾何学』，岩波書店， 1992.\par\smallskip
\noindent [357]\ Taubes, C., Self-dual connections on non-self-dual manifolds, J. Diff. Geom. 17(1982), pp.139-170.\par\smallskip
\noindent [358]\ Taubes, C., Casson's invariant and gauge theory, J . Diff. Geom. 31(1990), pp. 547-599.\par\smallskip
\noindent [359]\ Taubes, C., SW • Gr: from the Seiberg-Witten equations to pseudoholomorphic curves, J . Amer. Math. Soc. 9(1996), pp. 845-918.\par\smallskip
\noindent [360]\ Taubes, C., The structure of pseudo-holomorphic subvarieties for a degenerate almost complex structure and symplectic form on S1 x B叫 Geom. Topol. 2(1998), pp. 221-332.\par\smallskip
\noindent [361]\ Taubes, C., Moduli spaces and Fredholm theory for pseudoholomorphic subva— rieties associated to self-dual harmonic 2-forms, preprint.\par\smallskip
\noindent [362]\ Thom, R., Sur une partition en cellules associee a une fonction sur une variete, C. R. Acad. Sci. Paris 228(1949), pp. 973--975.\par\smallskip
\noindent [363]\ Thom邸， B., Eli邸hberg, Y. and Giroux, E., "3-dimensional contact geometry'' in Contact and Symplectic Geome乃，Thomas, C.(ed.), Cambridge Univ. Press, 1996.\par\smallskip
\noindent [364]\ Thomas, R., A holomorphic Casson invariant for Calabi-Yau 3-folds, and bundles on K3 fibrations. J. Differential Geom. 54(2000), no. 2, 367-438. math/9806111.\par\smallskip
\noindent [365]\ Tian, G., "Smoothness of the universal deformation space of compact Calabi-Yau manifolds and its Petersson-Weil metric" in Mathematical Aspects of String Theory, World Scientific, Singapore, 1987, pp. 543-559.\par\smallskip
\noindent [366]\ Tian, G., Gauge theory and calibrated geometry I , preprint.\par\smallskip
\noindent [367]\ Todorov, A., The Weil-Petersson geometry of the moduli space of SU(n;;: 3) (Calabi-Yau) manifolds, I . Comm. Math. Phys. 126(1989), pp. 325-346.\par\smallskip
\noindent [368]\ 戸田宏·三村護『ホモトピー論』（紀伊國屋数学叢書3), 紀伊國屋書店， 1975-\par\smallskip
\noindent [369]\ 戸田宏·三村護『リー群の位相上·下』（紀伊國屋数学叢書14A, B), 紀伊國屋書 店. 1979.\par\smallskip
\noindent [370]\ 上野健爾，『代数幾何』，岩波書店， 2005.\par\smallskip
\noindent [371]\ 上野健爾•清水勇二，『複素構造の変形と周期 共形場理論への応用\textasciitilde{} 岩 波書店， 2008.\par\smallskip
\noindent [372]\ Uhlenbeck, K. and Yau, S. T., On the existence of Hermitian-Yang-Mills connections in stable vector bundles, Comm. Pure. Appl. Math. 39(1986), pp. 257-293.\par\smallskip
\noindent [373]\ Vafa, C., "Topological mirrors and quantum rings" in Essays on Mirror Manzfolds, Yau, S. T.(ed.), Internatio叫Press, Hong-Kong, 1992, pp. 96-119.\par\smallskip
\noindent [374]\ Vafa, C., Extending Mirror conjecture to Calabi-Yau with bundles. hepth/9804031.\par\smallskip
\noindent [375]\ Vanderbauwhede, A., "Center manifolds, normal forms and elementary bifurcations" in Dynamics Reported, Kirchgraber, U. and Walther, H.(eds.), John Wiley, 1989.\par\smallskip
\noindent [376]\ Vassiliyev, V., Lagrange and Legendre Characteristic Classes, Advanced Studies in Contemporary Math., Gordon and Breach, 1988.\par\smallskip
\noindent [377]\ Vinogradov, A., The Q . ' . spectral sequence, Lagrangian formalism and conservation laws I , II, J. Math. Anal. Appl. 100(1984), pp.1-40, 41-129.\par\smallskip
\noindent [378]\ Viterbo, C., A proof of Weistein conjecture in JR叫Ann. Inst. H. Poincare, Anal. Non Lineaire 4(1987), pp. 337-357.\par\smallskip
\noindent [379]\ Viterbo, C., A new obstruction to embedding Lagrangian tori, Invent. Math. 100(1990), pp. 301-320.\par\smallskip
\noindent [380]\ Viterbo, C., Symplectic topology as the geometry of generating functions, Math. Annal. 292(1992), pp. 685-710.\par\smallskip
\noindent [381]\ Viterbo, C., Exact Lagrange submanifolds, periodic orbits and the cohomology of free loop spaces, J . Di.ff. Geom. 47(1997), pp. 420-468.\par\smallskip
\noindent [382]\ Wang, G. and Ye, R., Equivariant and Bott-type Seiberg-Witten Floer Homology: Part I . math/9901058.\par\smallskip
\noindent [383]\ Weil, A., Elliptic Functions According to Eisenstein and Kronecker, Sprmger, Berlin, 1976.\par\smallskip
\noindent [384]\ Weinstein, A., Symplectic manifolds and their Lagrangian submanifolds, Adv. in Math. 6(1971), pp. 202-213.\par\smallskip
\noindent [385]\ Weinstein, A., Periodic orbits for convex Hamiltonian system, Ann. of Math. 108(1978), pp. 507-518.\par\smallskip
\noindent [386]\ Weinstein, A., "Noncommutative geometry and geometric quantization" in Symplectic Geometry and Mathematical Physics, Donato, P. et al.(eds.), Birkhauser, Boston, 1991, pp.11-15.\par\smallskip
\noindent [387]\ Weinstein, A., Classical theta functions and quantum tori, Publ. RIMS 2(1994), pp. 327-333.\par\smallskip
\noindent [388]\ Weinstein, A., Deformation quantization, Seminaire Bourbaki 1993/94, No. 789, Asterisque 227, 1995, pp. 389-409.\par\smallskip
\noindent [389]\ Weitsman, J., Real polarization of the moduli space of flat connections on a Riemann surface, Comm. Math. Phys. 145(1992), pp. 423-433.\par\smallskip
\noindent [390]\ Whitney, H., The self intersections of a smooth n-manifold in 2n-space, Ann. of Math. 45(1944), pp. 220-246.\par\smallskip
\noindent [391]\ Witten, E., Supersymmetry and Morse theory, J. Diff. Geom. 117(1982), pp. 353-386.\par\smallskip
\noindent [392]\ Witten, E., Topological sigma model, Comm. Math. Phys. 118(1988), pp. 411-449.\par\smallskip
\noindent [393]\ Witten, E., Quantum field theory and Jones polynomial, Comm. Math. Phys. 121(1989), pp. 661-692.\par\smallskip
\noindent [394]\ Witten, E., Two dimensional gauge theory revisited, J. Geom. Phys. 9(1992), pp. 303-368.\par\smallskip
\noindent [395]\ Witten, E., "Chern-Simons gauge theory as a string theory" in The Floer Memorial Volume, Hofer, H. et al.(eds.), Birkh釦user, Basel, 1995, pp. 637-678.\par\smallskip
\noindent [396]\ Witten, E., "Phases of N = 2 theories in two dimensions" in Mirror Symmetry II, Yau, S. T.(ed.), International Press, Hong-Kong, 1996.\par\smallskip
\noindent [397]\ Witten, E., "Mirror manifolds and topological field theory" in Mirror Symmetry I , Yau, S. T.(ed.), International Press, Hong-Kong, 1998.\par\smallskip
\noindent [398]\ Witten, E. and Olive, 0., Super Symmetric algebra that include topological charg, Physics Lett. 70B(1978), pp. 97-101.\par\smallskip
\noindent [399]\ Woodhause, N., Geometric Quantization, Oxford Univ. Press, 1992.\par\smallskip
\noindent [400]\ Yau, S. T., On the Ricci curvature of a compact Kahler manifold and the complex Monge-Ampらre equation I , Comm. Pure Appl. Math. 31(1978), pp. 339-411.\par\smallskip
\noindent [401]\ Yau, S. T.(ed.), Mirror Symmetry II, AMS/IP Studies, International Press, Hong-Kong, 1997.\par\smallskip
\noindent [402]\ Yau, S. T.(ed.), Mirror Symmetry I , AMS/IP Studies, International Press, HongKong, 1998.\par\smallskip
\noindent [403]\ Yau, S. T.(ed.), Mi汀or Symmetry III, AMS/IP Studies, International Press, Hong-Kong, 1999.\par\smallskip
\noindent [404]\ Ye, R., Gromov's compactness theorem for pseudo holomorphic curves, Trans. Amer. Math. Soc. 342(1994), pp. 671-694.\par\smallskip
}